\subsubsection{紧算子的谱}\label{subsec:compact-operator-spectrum}

本节证明关于紧算子谱的若干结果. 由于紧算子有界, 紧算子的谱具有 Section~\ref{subsec:bounded-operator-spectra} 中所述的全部性质. 当然, 加上紧性假设后, 可以得到更多结论. 这里只限于 Hilbert 空间上的算子, 不过所给出的许多结果都可以推广到 Banach 空间上的算子.

在 Hilbert 空间中, 可以利用投影定理及其推论. 为利用这一点, 先描述有限秩算子的谱.

\begin{Lemma}
\label{lem:finite-rank-spectrum-alternative}
假设 $A\in\mathcal{L}(H)$ 具有有限秩. 那么对每个 $\lambda\in\mathbb{C}\backslash\{0\}$, 以下两种情形恰有一种成立:\begin{enumerate}\item $\lambda\in\rho(A)$, 或者\item $\lambda\in\sigma_p(A)$. 此时 $\lambda$ 为有限重数的特征值.\end{enumerate}
\end{Lemma}

证明直接来自线性代数中的相应结果, 留给读者 (Problem~\ref{ex:finite-rank-spectrum-alternative-proof}).

现在对有限秩算子证明 Fredholm 二择一定理的一个略有不同的版本. 这一版本实际上只是一个技术性结果, 在证明下面的解析 Fredholm 定理 (Theorem~\ref{thm:analytic-fredholm}) 时将会用到.

\begin{Lemma}
\label{lem:analytic-finite-rank-alternative}
设 $G\subset\mathbb{C}$ 为区域, 并假设
\begin{equation}\label{eq:analytic-finite-rank-operator-family}
\mathbb{C}\supset G\ni\lambda\mapsto F(\lambda)\in\mathcal{L}(H)
\end{equation}
在 $G$ 中解析. 进一步假设, 对每个 $\lambda\in G$, $F(\lambda)$ 均具有有限秩且
\begin{equation}\label{eq:analytic-finite-rank-fixed-range-subspace}
\mathcal{R}(F(\lambda))\subseteq M,
\end{equation}
其中 $M$ 为 $H$ 的一个不依赖于 $\lambda$ 的有限维子空间. 那么以下两者之一成立:\begin{enumerate}\item 对任何 $\lambda\in G$, $(I-F(\lambda))^{-1}$ 均不存在, 或者\item 对每个 $\lambda\in G\backslash S$, $(I-F(\lambda))^{-1}$ 均存在, 其中 $S$ 是 $G$ 中的离散集 (即在 $G$ 中没有极限点). 此时函数 $\lambda\mapsto(I-F(\lambda))^{-1}$ 在 $G\backslash S$ 上解析, 且若 $\lambda\in S$, 则 $F(\lambda)\phi=\phi$ 有有限维解族.\end{enumerate}
\end{Lemma}

\begingroup
\postdisplaypenalty=10000
\renewcommand{\endProof}{\par\nobreak\hfill$\blacksquare$\par\vspace{0.1cm}}
\begin{Proof}
设 $\{\psi_i\}_{i=1}^N$ 为 $M$ 的一组基. 那么存在解析向量函数
\begin{equation}\label{eq:analytic-finite-rank-vector-functions}
G\ni\lambda\mapsto\gamma_i(\lambda)\in H,\quad i=1,\ldots,N,
\end{equation}
使得
\begin{equation}\label{eq:analytic-finite-rank-family-expansion}
F(\lambda)\phi=\sum_{i=1}^n(\gamma_i(\lambda),\phi)\psi_i.
\end{equation}
设 $\Lambda(\lambda)$ 为分量如下的 $N\times N$ 矩阵:
\begin{equation}\label{eq:analytic-finite-rank-matrix-components}
\Lambda_{ij}(\lambda)=(\gamma_j(\lambda),\psi_i).
\end{equation}
请读者验证, $F(\lambda)\phi=\phi$ 有非平凡解, 当且仅当
\begin{equation}\label{eq:analytic-finite-rank-invertibility-determinant}
d(\lambda):=\det(I-\Lambda(\lambda))=0.
\end{equation}
然而, $d(\lambda)$ 在 $G$ 上解析. 因而由复变函数的一个标准结果, 要么 $d$ 在 $G$ 中恒为零, 要么 $d$ 的零点形成一个离散集.

由于 $F$ 的值域有限维, $F(\lambda)\phi=\phi$ 的解空间也有限维. 证明完毕.
\end{Proof}
\endgroup
